\section{Experimental Methodology}
\label{sec:methodology}

The methodology is specified before implementation. The companion experimental specification supplies implementation-ready contracts; this section fixes the scientific comparison. Feasibility pilots resolve testbed-dependent quantities but are excluded from all confirmatory estimates.

\subsection{Research questions, hypotheses, and sole primary estimand}

For each scheduled policy attempt in randomized block $j$, let $A_{bj}=1$ only when policy $b\in\{B_0,B_1,B_2,B_3\}$ passes every applicable mandatory performance-blind per-attempt acceptance gate; otherwise $A_{bj}=0$. Let $E_j^{\mathrm{mem}}=1$ only when the final B0/B1 page record is \textsc{Pair-Equivalent}. Define pair indicators $A_j^{(1,0)}=A_{B_1j}A_{B_0j}E_j^{\mathrm{mem}}$, $A_j^{(2,1)}=A_{B_2j}A_{B_1j}$, and $A_j^{(3,2)}=A_{B_3j}A_{B_2j}$. The run-level sample p99.9 external turnaround latency $Y_{bj\ell}$ enters a contrast only when the corresponding pair indicator is one. RQ1 estimates CPU affinity ($B_1$ versus $B_0$); RQ2 the incremental effect of memory binding ($B_2$ versus $B_1$); and RQ3 the incremental effect of the complete mapping ($B_3$ versus $B_2$). RQ4 treats packet size, offered load, traffic shape, flow count, queue count, and worker count as effect modifiers. RQ5 evaluates cost-term ablations and explanatory counters. RQ6 evaluates policy-rank stability under a declared workload or topology change.

There is one confirmatory estimand. A single reference condition $\ell^\star$ is frozen after a performance-blind feasibility pilot using the decision rules in the campaign protocol. Let $\set J_{32}=\{j:A_j^{(3,2)}=1\}$ among the blocks scheduled at $\ell^\star$. For $j\in\set J_{32}$,
\begin{equation}
 D_j=\log\!\left(\frac{Y_{B_3j\ell^\star}}{Y_{B_2j\ell^\star}}\right),
 \qquad
 \theta_{32}^{\mathrm{acc}}=\operatorname{E}\!\left[D_j\mid
 A_j^{(3,2)}=1,\ j\text{ scheduled at }\ell^\star\right].
 \label{eq:primary-estimand}
\end{equation}
The expectation is over accepted independent pairs under the frozen testbed and randomized run-order procedure; $\exp(\theta_{32}^{\mathrm{acc}})-1$ is the conditional geometric relative change in p99.9. It is not an unconditional effect over every scheduled block. The directional hypotheses are $H_0:\theta_{32}^{\mathrm{acc}}\geq0$ and $H_1:\theta_{32}^{\mathrm{acc}}<0$. The confirmatory decision uses a one-sided paired $t$-test at $\alpha=0.05$ on the independent accepted-pair values $D_j$, with the aligned one-sided upper confidence bound for $\theta_{32}^{\mathrm{acc}}$; a two-sided 95\% Student interval is also reported for estimation. RQ1 and RQ2 use the analogous accepted-pair sets defined by $A_j^{(1,0)}$ and $A_j^{(2,1)}$. Their effects, other workloads, absolute differences, and all p99.99 results are key secondary or exploratory outcomes and cannot substitute for \cref{eq:primary-estimand}.

No superiority claim is made before this test. A latency change accompanied by a different packet path, insufficient goodput, or a failed correctness gate is not evidence for the hypothesis. Operational feasibility is a mandatory secondary outcome: each policy's attempt-acceptance rate and each adjacent pair's completion rate are reported without latency-based selection. A policy is not described as practically beneficial when its latency estimate comes from a small selected subset and it has materially greater planning, application, or validation attrition.

\subsection{Required testbed and measurement chain}

Confirmatory measurement requires two physical hosts: an isolated traffic-generator/measurement host (TG) and the multi-socket SUT. A one-host loopback or colocated-generator configuration is diagnostic only because it omits or perturbs the physical links and competes for SUT resources. Direct links are preferred; any switch, cable, or optical path is identified, isolated, and invariant.

The primary observation follows this complete chain:
\begin{quote}\small
TG request timestamp and transmit $\rightarrow$ physical link $\rightarrow$ SUT NIC ingress $\rightarrow$ RX descriptor and payload DMA $\rightarrow$ poll/IRQ progress $\rightarrow$ optional event-ring handoff $\rightarrow$ worker and application state $\rightarrow$ response construction $\rightarrow$ TX payload/descriptor publication and doorbell $\rightarrow$ NIC transmit $\rightarrow$ physical link $\rightarrow$ TG receive timestamp.
\end{quote}
Both primary timestamps are in one TG clock domain, so no inter-host clock subtraction is required. The response is called external wire-to-wire \emph{turnaround} latency, not one-way SUT latency. Constant link and TG contributions are not subtracted from the primary comparison.

Timestamp capability is selected by a frozen decision tree. First preference is an external tap/FPGA or equivalent instrument that timestamps request egress and response ingress in one validated hardware clock. Second is correlated TX/RX hardware timestamps from the TG NIC's same physical hardware clock, provided matching capacity, rate behavior, timestamp loss, and capture-layer semantics pass validation. Third is TG software timestamps; these are diagnostic unless their end-to-end error bound, loss rate, and perturbation are established below a predeclared fraction of the smallest relevant effect. One-way SUT timestamps are secondary and require validated physical-clock synchronization, pre/post offset and drift bounds, and a timestamp layer stated precisely. Merely reporting PTP lock is insufficient.

The clock record contains source, clock identifier, capture boundary, nominal resolution, empirical quantization, conversion, wrap width and extension rule, calibration interval, timestamp loss, and overhead. Pre/post loop or known-delay trials bound bias and drift. A run fails if timestamp ambiguity or the error bound exceeds its frozen gate.

\subsection{Environment and invariants}

The environment table reports CPU model/stepping/microcode, sockets, dies, cores, PUs, SMT, caches, NUMA nodes/distances/memory population, socket links, firmware and power settings, kernel/boot/isolation configuration, NIC/ports/firmware/driver/BDF/full exposed ancestry, PCIe link state, IOMMU and DDIO state where observable, queues/descriptors/offloads/MTU, huge pages, DPDK or AF\_XDP versions/modes, compiler/linker and binary digest, TG and timestamp hardware/software, and counter availability. No NIC model, rate, or counter is assumed before stand qualification.

Within a matched block, the following are identical: source and binary; compiler/linker flags and link order; application logic; frontend profile; NIC firmware/driver/link; queue count and fixed $(r_q,t_q)$ pairs; RSS key, table, hash fields, and flow-to-queue rule; ownership cardinalities and frozen catalog digest; descriptors, offloads, MTU, burst/batch limits; packet format and response semantics; pool/ring usable capacities and page size; worker-role, I/O-role, and physical-execution-context inventories and counts; RTC/staged grouping rules; safety cpuset; frequency/power/isolation settings; TG trace and exact expected sequence set; timestamp method; phase lengths; telemetry; and acceptance gates. $B_0$--$B_2$ also share the topology-agnostic reference mapping $\set R_0$; $B_3$ may change only the controllable finite logical mappings explicitly included in its treatment. Process lifetime, allocation/fault order, and reset procedure are fixed to prevent a later policy from inheriting another policy's pages or queue state \cite{mytkowicz2009wrongdata}.

The DPDK profile freezes driver binding, environment abstraction layer parameters, RX/TX setup, mempool and descriptor construction, cache/burst settings, polling/completion mode, and self-reported owners. The AF\_XDP profile freezes native/generic and copy/zero-copy modes, UMEM frame layout, RX/TX/fill/completion rings, need-wakeup, busy-poll/interrupt behavior, XDP program, and socket-to-queue binding. Confirmatory comparisons use one frozen primary profile with single-owner queue progress whenever a controllable role exists. A multi-owner profile requires a separate finite ownership catalog and share-validation capability record and remains secondary or a replication until that gate passes; a second frontend profile is likewise a replication, not another baseline.

\subsection{Topology-blind initialization and page-equivalence gate}

The B0/B1 contrast uses a separate initialization context $x_{\mathrm{init}}$ that never participates in the measured packet-processing path. In both policies it runs on the same CPU selected from the initialization cpuset by the smallest stable operating-system CPU identifier after functional exclusions. This rule reads no NIC BDF, PCIe ancestry, NUMA identity or distance, page location, cost term, or latency outcome; it is an experimental invariant, not a locality optimization. Under the ordinary OS memory policy it creates fresh objects in the same stable type/identity order and faults every required page with the same access pattern. Faulting completes before any data-plane physical-execution-context affinity is applied. The process, allocations, and frontend state are fresh for every attempted run, so no policy inherits pages or warmed queue state.

For baseline $b$, run $j$, and material object $m$, let $\widetilde\nu_{bjm}(n)$ be the independently observed fraction of covered pages on node $n$. Desired distribution, requested memory policy, and $\widetilde\nu$ are separate records. Each attempt first emits an immutable page-initialization record and independently validates the coverage and internal consistency of its own observed vector. It does not compute a pair distance when the matching policy record is absent; the RQ1 relation is then \textsc{Pair-Pending}, and the attempt may continue through measurement under the fixed randomized schedule. After both matching B0 and B1 records exist, the validator or result aggregator joins mandatory objects and evaluates
\begin{equation}
 d_{\mathrm{TV}}(\widetilde\nu_{B_0jm},\widetilde\nu_{B_1jm})
 =\frac{1}{2}\sum_{n\in\set N}
 \left|\widetilde\nu_{B_0jm}(n)-\widetilde\nu_{B_1jm}(n)\right|
 \leq\tau_{\mathrm{mem}}
 \label{eq:page-equivalence}
\end{equation}
for every mandatory material payload, controllable descriptor, event-ring, and state object, with compatible observation coverage. A performance-blind pilot selects $\tau_{\mathrm{mem}}$ from the page-location method's resolution and repeatability under repeated executions of this invariant initializer; the rule and resulting value are frozen before confirmation and cannot be revised after latency is visible. The final pair status is \textsc{Pair-Equivalent}, \textsc{Page-Equivalence-Unsupported}, or \textsc{Page-Equivalence-Failed}. The latter two cannot enter RQ1, but they do not retroactively invalidate an otherwise accepted attempt for another pre-specified contrast. Pair evaluation uses only immutable page records, never latency; it causes no rerun, substitution, schedule delay, or change to randomized order. Passing this gate establishes measured equivalence within $\tau_{\mathrm{mem}}$, not equality of unobserved physical pages.

\subsection{Strictly incremental placement policies}

Hardware queues do not have CPU masks in the abstract. Their CPU association is induced by the poller, IRQ, or completion role that owns them and by the physical execution context realizing that role. Resource inventories, counts, fixed queue pairs, flow steering, ownership cardinalities, and finite ownership catalogs are common invariants. The topology-agnostic $\set R_0$ fixes queue-to-worker identities and catalog indices for $B_0$--$B_2$. $B_3$ may optimize those logical relations only over controllable catalog entries while also selecting physical execution-context PUs and memory; an observed-only share remains its frozen workload-dependent proxy and is not a B3 decision. This ownership change is a declared treatment component. In the primary single-owner profile, every non-null direction has $k_q^d=1$ and share one in all four policies.

\begin{table}[t]
\centering
\small
\begin{tabularx}{\linewidth}{@{}p{0.19\linewidth}YY@{}}
\toprule
Policy & Exactly controlled & Explicitly unchanged or delegated \\
\midrule
$B_0$ default & Functional launch inside the common safety cpuset, fixed $\set R_0$, and the invariant topology-blind initializer & No experiment-selected data-plane affinity or NUMA policy; fixed queues, RSS, catalogs, RTC/staged semantics, and all other invariants remain common \\
$B_1$ CPU only & Exact B0 initialization and page-equivalence gate, plus deterministic affinity for every data-plane physical context & No memory-policy or first-touch change; no change to $\set R_0$, queues, RSS, pools, descriptors, rings, state, or sizes \\
$B_2$ CPU + memory & Exact B1 execution-context assignments and $\set R_0$, plus the CPU-access-share memory rule on fresh allocations & No PCIe-root/path term, ownership-catalog change, load-aware joint re-planning, or non-memory change \\
$B_3$ complete & Fresh allocation plus the complete \textsc{Feasible} placement, including a controllable finite ownership entry, occupancy/PUs, and location-qualified objects & Queue identities, RSS, workload, RTC/staged semantics, catalogs, and non-placement invariants remain fixed; fallback or limited-search output is not pooled \\
\bottomrule
\end{tabularx}
\caption{Strictly incremental policy definitions. A baseline changes only the control added in its row.}
\label{tab:baselines}
\end{table}

$B_0$ is reproducible rather than “whatever happened”: its allowed masks, actual CPUs, default page policy, topology-blind initializer, observed page distributions, queue counters, and ownership observations are archived. If a frontend cannot operate without explicit CPU assignment, its unmodified documented launch mapping is the $B_0$ default and is reported as such; the study does not call this an operating-system choice.

$B_1$ orders eligible physical cores by stable OS identifier and assigns active data-plane physical execution contexts round-robin after reserved cores, respecting only functional capacity, SMT, and isolation constraints. Each context is assigned once even when it realizes both RTC poller and worker roles; distinct staged contexts remain distinct occupants. It does not inspect NIC BDF, PCIe ancestry, NUMA distance, page location, or calibrated path cost. This affinity is applied only after the common B0/B1 initializer has completed, and the RQ1 pair is admitted only after the pair-level evaluation in \cref{eq:page-equivalence} returns \textsc{Pair-Equivalent}. Consequently B1 differs from its paired B0 attempt only by data-plane CPU affinity.

$B_2$ preserves every $B_1$ CPU assignment and the complete reference mapping $\set R_0$. For each material memory object $m$, the frozen pipeline contract must list a nonempty set of its actual CPU accessors $\set X_m^{\mathrm{cpu}}$ and nonnegative demand shares $\sigma_{mx}$ summing to one. These may include progress contexts, payload workers, both event-ring endpoints, state users, or another context justified by the frozen packet path. A material object without a justified CPU accessor makes the requested $B_2$ profile undefined and must be resolved before freezing the campaign; NIC locality is not used as a substitute. With $N(c)$ the NUMA node containing PU $c$, $B_2$ selects
\begin{equation}
 n_m^{B_2}=\arg\max_{n\in\set N}
 \sum_{x\in\set X_m^{\mathrm{cpu}}}\sigma_{mx}
 \mathbf{1}[N(\eta_{B_1}(x))=n],
 \label{eq:b2-memory-rule}
\end{equation}
using the smallest stable node identifier on a tie. It creates fresh allocations, records desired distribution and requested policy separately, binds or faults $m$ according to the frozen per-object mechanism before applying the preserved B1 data-plane affinity, and verifies the observed distribution. Shares come from operation counts or performance-blind CPU-access calibration, never NIC/PCIe locality or baseline latency. A secondary sensitivity contrasts this aggregate-node rule with a predeclared dominant-accessor rule and demand-share-weighted interleaving when multiple nodes carry material CPU demand; it does not redefine primary $B_2$. Thus $B_2$ equals exact $B_1$ CPU placement plus CPU-relative memory binding and nothing else.

$B_3$ alone uses the complete topology/load score and may replace $\set R_0$ with a feasible optimized $\set R=(\omega,\kappa^{\mathrm{rx}},\kappa^{\mathrm{tx}})$. Three pre-specified secondary mechanism analyses separate ownership and physical placement: full $B_3$; $B_3$ with $\set R=\set R_0$ while physical placement is optimized; and controllable ownership/worker mapping optimized while every other assignment is frozen at its verified $B_2$ value. The third contrast is executed only if all alternatives remain feasible and preserve packet semantics and offered work. A shared-owner ablation additionally compares frozen finite catalog entries or, for observed-only sharing, strata defined by the frozen observed-share acceptance range; it never estimates a continuous share choice or claims that an observed share was applied. These analyses and the separate topology-only versus load-aware contrast are not additional levels in the four-baseline ladder.

\subsection{Workload matrix and run protocol}

Packet sizes include the minimum supported Ethernet frame, justified interior sizes, the standard-MTU boundary, and a jumbo boundary only if the invariant profile enables it. On-wire accounting states whether preamble, inter-packet gap, and frame check sequence are generated/reported. Queue and worker counts span one through the largest value jointly supported by NIC queues and the isolated PU set using a reduced power-of-two design plus the boundary value; feasible $Q<W$, $Q=W$, and $Q>W$ cases exercise sharing/conflict behavior. Flow counts include one, at least one independently steerable flow per RX queue, and the smallest count whose repeated RSS probe meets the frozen occupancy-stability rule.

Absolute rates are determined on the final stand by a feasibility campaign excluded from inference. A common zero-loss capacity is the minimum validated capacity across policies under the invariant path. Low, interior, first-knee, and highest common pre-gate rates are then selected by a pre-specified rate-sweep rule, not by visually favorable p99 differences. Constant-rate traces isolate placement; replayable ON/OFF traces vary burst factor and duration in units of calibrated service and queue-drain time. Uniform and skewed flow demands are included. The trace digest is identical within each block.

Each attempted run follows one auditable dependency order: (1) reset the environment and verify the neutral prior state; (2) create fresh allocations; (3) perform the invariant topology-blind allocation/fault sequence for B0/B1 or the policy-specific allocation/fault sequence for B2/B3; (4) collect and validate the attempt's own immutable page-location observation, leaving a B0/B1 contrast \textsc{Pair-Pending} if its mate is not yet available; (5) apply data-plane affinity and the catalog-defined ownership mapping; (6) complete independent read-back of desired, requested, and observed state; (7) execute fixed warm-up; and (8) measure the committed exact transmitted set, drain, post-validate, and cool down. When both matching page records later exist, the validator or aggregator performs the pair-level join and \cref{eq:page-equivalence} without changing this lifecycle. No policy inherits pages, process state, or warmed queues from another. If a frontend requires a different safe order, stand qualification must identify the unavoidable dependency and demonstrate with repeated performance-blind page-distribution observations that the alternate order does not let data-plane affinity alter B0/B1 memory placement; otherwise that frontend profile is unsupported for RQ1.

The pilot determines a conservative warm-up duration exceeding the longest observed stabilization time for throughput, frequency, temperature, and queue occupancy; that duration and tolerances are frozen. Measurement never stops when a confidence interval “looks precise.” Cooldown uses fixed time plus a pre-specified environmental gate; a failed gate consumes an attempted-run slot and is recorded.

Policy order is randomized within complete blocks defined by topology snapshot, frontend, trace, packet size, offered rate, queues/workers, and repetition. Runs from different boot/topology identities are not paired. Pilot runs, timestamp qualification, calibration, and confirmatory runs have disjoint identifiers and artifacts.

\subsection{Exact packet accounting and metrics}

Before transmission, the TG commits an arbitrary finite expected sequence set $E_j$ and its digest; contiguity is not required. For each $s\in E_j$, let $R_s$ be all received candidate responses carrying $s$, and let $R_s^{\mathrm{ok}}$ be those passing packet format, payload/integrity, request-response binding, and timestamp validity. The latency-eligible set is
\begin{equation}
 V_j=\{s\in E_j:|R_s|=|R_s^{\mathrm{ok}}|=1\},\qquad
 L_s=t_s^{\mathrm{rx}}-t_s^{\mathrm{tx}}\quad(s\in V_j).
 \label{eq:valid-sequence-set}
\end{equation}
The report separately counts missing sequences ($|R_s|=0$), corrupt/invalid responses ($|R_s|>0$ but $|R_s^{\mathrm{ok}}|=0$), duplicates ($|R_s|>1$), ambiguous matches, and out-of-range sequences. The no-valid-response rate is $1-|V_j|/|E_j|$; it is not labeled physical loss without the category decomposition. Goodput counts exactly one valid response per expected request. An initial prefix or terminal suffix cannot disappear because $E_j$ is known independently of receipt.

Let $N=|V_j|$ and let $L_{(1)}\leq\cdots\leq L_{(N)}$ be the one-indexed order statistics, retaining repeated values when timestamps tie. For $N>0$, the frozen nearest-rank estimator is
\begin{equation}
 \widehat Q_p=L_{(k_p)},\qquad
 k_p=\max\{1,\lceil pN\rceil\},\qquad 0\leq p\leq1.
 \label{eq:nearest-rank}
\end{equation}
Thus $p=0$ returns the minimum, $p=1$ the maximum, and every reported positive percentile uses $L_{(\lceil pN\rceil)}$ without interpolation. If $N=0$, every quantile is undefined and the run is invalid. Invalid, duplicate, corrupt, or ambiguous responses never enter $V_j$ and therefore cannot enter the order statistics. Raw tick differences are wrap-extended and converted once before stable numeric sorting; no intermediate rounding is allowed.

For a bounded-resolution histogram with ordered edges $e_0<\cdots<e_B$, bins are $[e_{b-1},e_b)$ except the final closed-right bin. The reported percentile is the unique bin interval whose cumulative count first reaches $k_p$; both interval bounds, $N$, and $k_p$ are returned rather than an interpolated point. Underflow or overflow containing $k_p$ gives an unbounded interval and is inadmissible for a primary percentile. The latency distribution reports the nearest-rank p50, p90, p99, and p99.9, maximum, interquartile range, median absolute deviation, and a pre-specified jitter statistic, with $N$ beside every percentile. p99.99 is always exploratory in this campaign. Throughput, goodput, missing/corrupt/duplicate/reorder rates, CPU utilization, cycles per valid response, queue occupancy/drops, planner and configuration time, page location, remote-memory indicators, cache events, memory bandwidth, and PCIe/inter-socket counters are secondary or explanatory.

Each run requires $|V_j|(1-0.999)\geq1000$, hence at least $10^6$ valid latency observations for p99.9. The TG therefore transmits a fixed $|E_j|$ chosen before confirmatory execution to satisfy this count under the frozen acceptance-loss allowance and the minimum sustained-time rule. Falling short invalidates the run; collection is not extended after inspecting latency. The count rule supports tail resolution but does not make packets independent.

Raw paired timestamps are the preferred artifact. If storage requires a histogram, its schema declares finite bin edges, interval closure, underflow/overflow counts, integer count width, total-count invariant, every nearest-rank index, the bin containing it, the maximum bin width around p99.9, and the resulting interval-censoring bound for every reported percentile. Such a histogram is described as bounded-resolution, not “lossless,” unless exact raw values are recoverable. Quantization and bin-width uncertainty are carried into a sensitivity report.

\subsection{Fixed sample size and statistical analysis}

An independent pilot estimates the standard deviation $s_D$ of paired log ratios and the performance-blind invalid-attempt fraction. Investigators specify a smallest scientifically relevant relative reduction $\delta_{\mathrm{rel}}$ from the latency budget before examining policy effects and convert it to $\delta^\star=|\log(1-\delta_{\mathrm{rel}})|$. For desired power $1-\beta_{\mathrm{II}}$ and one-sided $\alpha=0.05$, an initial normal-theory planning approximation is
\begin{equation}
 n=\left\lceil
 \left(\frac{z_{1-\alpha}+z_{1-\beta_{\mathrm{II}}}}{\delta^\star/s_D}\right)^2
 \right\rceil.
 \label{eq:power}
\end{equation}
The final planned $n$ is the smallest integer at or above this approximation whose one-sample noncentral-$t$ power for standardized effect $-\delta^\star/s_D$ reaches the target, after applying any predeclared conservative variance inflation. The protocol records the latency-budget rationale, pilot $s_D$, chosen power, approximation, noncentral-$t$ recomputation, rounding, and inflation. A fixed maximum number of attempted blocks is set by inflating $n$ for the pilot invalid fraction. The confirmatory campaign then runs exactly that attempt schedule; it does not add blocks based on an observed effect, $p$-value, or confidence width. If fewer than $n$ B3/B2 pairs pass performance-blind gates, the study reports reduced power and does not opportunistically continue. If $n_{\mathrm{acc}}<2$, the variance and $t$ statistic are undefined, so no confirmatory test is reported.

Let $n_{\mathrm{acc}}=|\set J_{32}|=\sum_j A_j^{(3,2)}$ be the number of accepted complete B3/B2 pairs from the fixed attempt schedule. The primary estimate is $\bar D=n_{\mathrm{acc}}^{-1}\sum_{j\in\set J_{32}}D_j$, reported with $\exp(\bar D)-1$. With sample standard deviation $s_D^{\mathrm{conf}}$, the primary test statistic and upper one-sided 95\% confidence bound are
\begin{equation}
 T=\frac{\bar D}{s_D^{\mathrm{conf}}/\sqrt{n_{\mathrm{acc}}}},\qquad
 U_{0.95}=\bar D+t_{0.95,n_{\mathrm{acc}}-1}
 \frac{s_D^{\mathrm{conf}}}{\sqrt{n_{\mathrm{acc}}}}.
 \label{eq:primary-ci}
\end{equation}
The one-sided paired $t$-test rejects $H_0$ only when $T<-t_{0.95,n_{\mathrm{acc}}-1}$, equivalently when $U_{0.95}<0$. The report also gives the two-sided 95\% Student interval $\bar D\pm t_{0.975,n_{\mathrm{acc}}-1}s_D^{\mathrm{conf}}/\sqrt{n_{\mathrm{acc}}}$. These quantities use run pairs, never individual packets.

The test assumes independent blocks in the target campaign, finite variance of $D_j$, and a sampling distribution of $\bar D$ adequately approximated by the Student model. Before unblinding, the analysis plan freezes diagnostics for block independence, a Q--Q display, skewness and tail-weight summaries, maximum studentized residual, and leave-one-pair-out influence. Diagnostics never authorize deletion. If predeclared asymmetry, tail, or influence thresholds are exceeded, the paired $t$ result remains the sole confirmatory rule but is labeled assumption-sensitive; a studentized paired bootstrap-$t$ interval and a median/sign analysis are reported as robustness checks. No alternative is selected because it is more favorable.

Swapping $B_2/B_3$ labels within pairs tests the sharp null of no pair-specific treatment effect, not the weak conditional mean null $H_0:\theta_{32}^{\mathrm{acc}}\geq0$. That label-swapping randomization test is therefore a secondary sensitivity analysis, using exact enumeration when feasible or a frozen-seed Monte Carlo approximation with reported simulation error. A block-resampled mean interval is an additional sensitivity only when $n_{\mathrm{acc}}\geq30$; no bias-corrected bootstrap is used with few blocks.

Within-run uncertainty of each sample p99.9 is reported separately using moving-block resampling or stationary subsampling over the time-ordered valid sequence stream, with block length selected from a pilot autocorrelation rule. Packets are never treated as independent experimental repetitions. RQ1 and RQ2 paired log-ratio tests form one key-secondary family controlled by Holm's procedure; workload interactions, ablations, RQ5, and RQ6 are labeled secondary/exploratory with their own declared families. Reports give planned and attempted block counts; accepted attempts separately for B0--B3; complete-pair counts for B1/B0, B2/B1, and B3/B2; and policy-specific gate-failure counts and rates. Failure categories distinguish planner, calibration/model domain, application, rollback, placement read-back, page verification, correctness, timestamp, telemetry, and environment, and identify which policy made a block incomplete. Planner/model reporting separately gives the proportions of \textsc{Search-Limited}, \textsc{Cap-Limited}, \textsc{Data-Incomplete}, \textsc{Calibration-Failed}, and application failure using explicit denominators; a campaign-wide calibration failure is one model-admission failure, not duplicated as a latency observation for every block.

No run is removed because its latency is inconvenient. Exclusions are determined without consulting $Y_{bj\ell}$ and are limited to the gates below, a recorded external stand fault, or a missing artifact designated mandatory in advance. All attempts and reasons remain in the index. There are no outcome-dependent reruns, neighboring-run substitutions, or arbitrary latency imputations for failed attempts. The complete-pair primary analysis remains unchanged. Differential attrition is an interpretation limit and a secondary operational-feasibility outcome. Weighting or bounding analyses are exploratory only when their covariates, positivity checks, formulas, and bounds were frozen before unblinding; they never replace the primary test.

\subsection{Correctness gates, validity, and reproducibility}

A run is admitted only if it passes packet format and integrity; exact $E_j$ accounting; the frozen missing/corrupt/duplicate/reorder thresholds; ring overwrite/lapping detection; producer/consumer progress and complete drain; logical-role-to-physical-context grouping and effective affinity; one-slot RTC and distinct-slot staged occupancy; its own desired/requested/observed location records and page-observation coverage for every material payload, descriptor, ring, and state object; reference or B3-selected finite queue-ownership entry, worker mapping, observed-share acceptance predicate, and RSS mapping; topology digest; clock/timestamp integrity; and mandatory telemetry. The later B0/B1 object-wise pair gate controls RQ1 pair admission, not the already frozen acceptance decision for an otherwise valid individual attempt. Unobservable mandatory page location yields the predeclared unsupported status rather than an assumed pass. A saturation study may intentionally violate a loss threshold but is exploratory and cannot enter \cref{eq:primary-estimand}. An attempt is rejected when one of its own gates fails; post hoc ``good intervals'' are forbidden.

Cost ablations re-plan after removing the owned component or group. Sensitivity covers component and group-level coefficients, topology-only versus load-aware planning, the three ownership/physical-placement mechanism variants, primary single-owner versus qualified finite shared-owner profiles, packet size/rate/burstiness/flow skew, queue/worker counts, SMT, local/cross-socket candidates, normalization, missing information, and one controlled topology change if available. Algorithm quality is evaluated against the independently complete exact-small and locality-only comparators specified in \cref{sec:algorithm}, including missed-feasibility rate; \textsc{Exact-Budget-Exceeded} cases have no certified optimum and do not change the main heuristic's status semantics.

Internal-validity threats include thermal/frequency drift, allocation history, RSS hashing, background work, interrupt leakage, order effects, and policy-dependent planning/application/read-back failure; isolation, fixed lifecycles, monitoring, randomized order, immutable gates, and policy-specific attrition reporting mitigate but do not remove them. Construct validity is limited by turnaround including invariant TG/link time and by imperfect counters; the explicit measurement chain and one clock domain bound interpretation. Conclusion validity is limited by dependent observations, extreme-quantile uncertainty, invalid runs, multiplicity, and selection into accepted pairs. Run-level pairing, fixed sample size, within-run uncertainty, pre-specified families, and operational-feasibility reporting address uncertainty, while the estimand remains explicitly conditional and differential attrition limits generalization. External validity remains bounded by the tested CPU, NIC, firmware, topology, and frontend.

The reproducibility package will contain topology/raw-source snapshots, environment and binary manifests, frontend capability and finite ownership catalogs, calibration and model artifacts including condition-coverage maps, candidate and planner audits, exact-small certificates and budget records, desired, requested, applied, and observed placements, immutable per-attempt page-initialization observations and pair-equivalence records, TG traces and exact $E_j$ sets, raw timestamps or bounded-resolution histograms with nearest-rank indices and containing intervals, clock validation, counters, gates, all planned, attempted, and accepted policy and pair indicators, categorized failures and denominators, the policy causing each incomplete block, exclusion reasons, the frozen attrition-sensitivity specification, the remaining analysis specification, and checksums. No measurement artifact or experimental result exists at the time of writing.
